\section{Optional optical-performance diagnostics}\label{sec:optional-optics}

These diagnostics concern a later optical-use problem. They do not replace the
two-face geometric maximum in \cref{sec:accuracy}; refractive-index modulation
is not part of the present shape-only objective.

\subsection{Optical objectives and a slope-sensitive identity}
Let $\mathcal P$ be a two-dimensional entrance pupil, $\bm\xi\in\mathcal P$ its
ray label, and $\bm X_\Gamma(\bm\xi)$ the ray's first intended intersection with
the interface. Every admitted ray must reach the intended transmitting patch
and have an unobstructed optical continuation. Let $\mathrm d\mu(\bm\xi)$ be a
nonnegative normalized illumination measure, $\int_{\mathcal P}\mathrm d\mu=1$.
Define $\overline\Psi=\int_{\mathcal P}\Psi(\bm X_\Gamma)\mathrm d\mu$ and
$W=\Psi(\bm X_\Gamma)-\overline\Psi$. With positive optical-path and angular
tolerances $W_{\rm ref}$ and $\theta_{\rm ref}$, respectively,
\begin{equation}
 J_W=\int_{\mathcal P}\frac{W^2}{W_{\rm ref}^2}\mathrm d\mu,
 \qquad J_d=\int_{\mathcal P}
 \frac{\norm{\Pg(\bm d_t-\bm d_i)}^2}{\theta_{\rm ref}^2}\mathrm d\mu.
 \label{eq:optical-objectives}
\end{equation}
Removing the optical piston does not move either conjugate. Tangential Snell
refraction gives the exact identity
\begin{equation}
 n_i\Pg(\bm d_t-\bm d_i)=\Gs\Psi.
 \label{eq:opd-slope-identity}
\end{equation}
A small-amplitude, rapidly varying path error can therefore give a large ray
error. Near grazing transmission, tangential control alone is poorly
conditioned for the full direction.

For a real image and a detector plane through $\bm I$ with unit normal
$\bm e_i$, require $\bm d_t\cdot\bm e_i>0$ and positive propagation distance.
Then $\bm Y=\bm X_\Gamma+[(\bm I-\bm X_\Gamma)\cdot\bm e_i]
\bm d_t/(\bm d_t\cdot\bm e_i)$. For a positive spot-radius tolerance,
\begin{equation}
 J_{\rm spot}=\int_{\mathcal P}\frac{|\bm Y-\bm I|^2}{r_{\rm ref}^2}\mathrm d\mu.
 \label{eq:spot-cost}
\end{equation}
Geometric ray errors are distinct from diffraction and measured optical quality.

\subsection{Fast optical modulation}
Let $\xi_a=\Rea(\Xi e^{-\mathrm i\omega t})$ be normal acoustic displacement,
and $\delta n_o,\delta n_i$ index perturbations along the ray segments
$\mathcal L_o,\mathcal L_i$. For real conjugates and a nominal Fermat ray,
\begin{equation}
 \delta\Psi=(n_o\bm d_o-n_i\bm d_i)\cdot\nn\,\xi_a
       +\int_{\mathcal L_o}\delta n_o\dd s
       +\int_{\mathcal L_i}\delta n_i\dd s.
 \label{eq:fast-optical-error}
\end{equation}
Here $s$ is arclength. Virtual conjugates use signed path segments. This is a
first-order estimate, not a complete acousto-optic scattering calculation.
Exposure time, index derivatives and the fast response determine its relevance.
Physical interface vibration still belongs to the geometric error even when
optical modulation is outside the experiment's objective.
